\documentclass[11pt,a4paper]{letter}
\usepackage[margin=1in]{geometry}
\usepackage[T1]{fontenc}
\usepackage[utf8]{inputenc}
\usepackage[hidelinks]{hyperref}

\signature{Jonathan Anthonio Nii Pedro Nelson\\Kwame Nkrumah University of Science and Technology (KNUST)\\Kumasi, Ghana\\\href{mailto:jonathannelson707@gmail.com}{jonathannelson707@gmail.com}\\ORCID: 0009-0006-9393-6237}
\address{Jonathan Anthonio Nii Pedro Nelson\\Kwame Nkrumah University of Science and Technology (KNUST)\\Kumasi, Ghana}
\date{\today}

\begin{document}
\begin{letter}{Prof. Keith Pilbeam\\Editor-in-Chief\\\textit{International Journal of Finance \& Economics}\\Wiley}
\opening{Dear Professor Pilbeam,}

I am pleased to submit the manuscript entitled ``Commodity Tail Dependence and Exchange-Rate Risk in Ghana: Copula Adequacy, Stress Tests, and Calendar Robustness'' for consideration in the \textit{International Journal of Finance \& Economics}.

The paper examines international commodity and exchange-rate risk in Ghana using daily-source data on cocoa, gold, Brent crude oil, WTI crude oil and the Ghanaian cedi over 2015--2026. Its central contribution is a robustness design that treats two defensible market-calendar constructions as co-equal specifications and asks which tail-dependence conclusions survive explicit marginal-model adequacy checks, absolute goodness-of-fit testing for eight static bivariate copulas, moving-block-bootstrap stress inference, and multiple-testing correction.

The findings have direct relevance to international finance and risk analysis. Gold--Brent, Gold--WTI and Brent--WTI reject all eight tested static copula families under both calendar constructions, whereas the three cocoa-related commodity pairs are calendar-sensitive. Although the COVID-19 and 2024 stress windows generate several nominally significant tail comparisons, none survives Bonferroni or Benjamini--Hochberg correction in either panel. Cedi-related copula and extreme-value results are reported cautiously because the tested continuous marginal models do not adequately reproduce the cedi's zero-heavy daily return process.

I believe the manuscript fits the journal's emphasis on exchange rates, financial uncertainty, risk management and econometric work with clear real-world relevance. The paper also provides a reproducible audit trail through a public project repository and distinguishes relative model selection from absolute model adequacy.

The research received no external funding, and the author declares no conflicts of interest. Substantive use of OpenAI's ChatGPT during manuscript preparation is disclosed transparently in the Methods section in accordance with Wiley's current guidance.

\textbf{AUTHOR CONFIRMATION BEFORE SUBMISSION:} Before uploading this letter, confirm that the manuscript has not been published previously and is not currently under consideration by another journal. Remove this sentence and replace it with the journal's required originality statement once confirmed.

Thank you for considering the manuscript.

\closing{Sincerely,}
\end{letter}
\end{document}
